\section{Introduction}
Small sets of objects (typically 1-3, sometimes up to 4) are enumerated rapidly and precisely, a phenomenon known as subitizing, while larger sets require estimation that follows Weber's law. One account explains this pattern through a single continuous magnitude code with scalar variability across the entire numerical range \citep{GallistelGelman2000}. Under this single-system view, small-number advantages reflect higher resolution at the low end of a continuous scale. An alternative account proposes two distinct representational systems: Parallel Individuation (PI), which tracks small sets of discrete items through ``object files'' and is often linked to subitizing, and the Approximate Number System (ANS), which represents larger quantities as noisy magnitude estimates with ratio-dependent discriminability \citep{FeigensonCarey2003,FeigensonDehaeneSpelke2004}.

These accounts make distinct predictions about neural similarity
structure. If numerosity is represented by a single continuous code, neural
similarity should scale with numerical ratio with no categorical
boundary. If PI and ANS are distinct, a categorical break should
separate their ranges at the PI capacity limit. The subitizing range
is often described as approximately 1-4 items \citep{TrickPylyshyn1994},
though reported limits vary across tasks and individuals
\citep{Cowan2001}, placing the
predicted boundary at either 3-4 or 4-5.

Neural similarity analysis can test which boundary model better captures
the representational structure. Representational similarity analysis
\citep[RSA;][]{Kriegeskorte2008} compares neural representational
dissimilarity matrices (RDMs) to model predictions. We apply
time-resolved RSA to single-trial EEG from 24 adults performing an
oddball change-detection task with 1-6 dot arrays. Using Linear
Discriminant Analysis (LDA) classifiers, which capture multivariate
spatiotemporal patterns from the EEG signal
\citep{KingDehaene2014}, we construct whole-epoch target RDMs
(6$\times$6) and 24-condition prime-target RDMs, comparing them against
PI models with boundaries at 3-4 and 4-5, a graded ANS model, and
reaction-time and pixel controls.

A complication in testing these predictions is the special status of
numerosity~1. Under Weber-law scaling, ``1'' is maximally distant from
all larger numerosities (the log-ratio between 1 and any other
numerosity is the largest possible), which can inflate ANS model fits. In addition,
numerosity~1 may be categorically privileged via a
singular-versus-plural (one-vs-more-than-one) distinction
\citep{Barner2007,Barner2008,Li2009}. To disentangle these effects, we
repeat all analyses after excluding pairs containing numerosity~1.

Across analyses, the data favor a dual-system account over a single
continuous magnitude code. ANS-like fits are strongest when numerosity~1
is included, but excluding ``1'' attenuates ANS fits while preserving PI
boundary structure. The 4-5 boundary outperforms the 3-4 alternative,
suggesting PI capacity at four items. Numerosity~1 is a distinct case
whose inclusion can increase the estimated fit of continuous models.
